\chapter{Dyskretna teoria Morse'a Formana}
\label{ch:tg-forman}

W klasycznej teorii Morse'a topologia rozmaitości jest
odczytywana z punktów krytycznych funkcji gładkiej.
W kompleksie symplicjalnym nie mamy pochodnych,
ale dysponujemy relacją bycia ścianą innego sympleksu.
Dyskretna teoria Morse'a wykorzystuje \emph{acykliczne
parowania komórek} do usuwania zbędnych par sympleksów.
Przedstawimy najpierw tę konstrukcję, a następnie
wyjaśnimy dokładnie, które wnioski da się udowodnić
samą algebrą kompleksów łańcuchowych.

\section{Relacja incydencji i parowanie komórek}
\label{sec:tg-forman-matching}

\begin{definicja}[Pokrycie w zbiorze częściowo uporządkowanym ścian]
Niech $K$ będzie skończonym kompleksem symplicjalnym.
Zgodnie z rozdziałem \ref{ch:tg-triangulacje} rozważamy
wyłącznie niepuste sympleksy.
Mówimy, że sympleks $\beta$ \emph{pokrywa} sympleks
$\alpha$, i piszemy $\alpha\prec\beta$, gdy
$\alpha\subsetneq\beta$ oraz
$\dim\beta=\dim\alpha+1$. \emph{Graf Hassego}
relacji ścian ma wierzchołki odpowiadające niepustym
sympleksom i krawędź między $\alpha$ i $\beta$
dokładnie wtedy, gdy $\alpha\prec\beta$.
\end{definicja}

\begin{definicja}[Dyskretne pole gradientowe]
\label{def:tg-discrete-morse-field}
Niech $K$ będzie skończonym kompleksem symplicjalnym.
\emph{Parowaniem} $\mathcal V$ nazywamy rodzinę par
$(\alpha^p,\beta^{p+1})$, w których
$\alpha\prec\beta$ i żaden sympleks nie występuje
w więcej niż jednej parze. Sympleks niewystępujący
w żadnej parze nazywamy \emph{krytycznym}.
\emph{$\mathcal V$-ścieżką wymiaru $p$} nazywamy ciąg
\[
\alpha_0^p\prec\beta_0^{p+1}\succ\alpha_1^p
\prec\cdots\prec\beta_{r-1}^{p+1}\succ\alpha_r^p,
\]
w którym $r\ge1$, dla $0\le i<r$ każda para
$(\alpha_i,\beta_i)$ należy do $\mathcal V$,
a $\alpha_{i+1}\prec\beta_i$ i
$\alpha_{i+1}\ne\alpha_i$. Parowanie jest
\emph{acykliczne}, gdy nie istnieje niebanalna
zamknięta $\mathcal V$-ścieżka z
$\alpha_r=\alpha_0$.
Acykliczne parowanie nazywamy dyskretnym polem
gradientowym.
\end{definicja}

To ostatnie ograniczenie jest kluczowe. Samo parowanie
może krążyć w zamkniętym cyklu i nie pozwalać na
monotoniczne upraszczanie kompleksu. Acykliczność
stanowi kombinatoryczny odpowiednik braku zamkniętych
trajektorii gradientowych.

\begin{stwierdzenie}[Charakteryzacja przez graf skierowany]
\label{thm:tg-morse-acyclic}
Niech $\mathcal V$ będzie parowaniem sympleksów
skończonego kompleksu $K$. W grafie Hassego skierujmy
każdą krawędź nieparowaną od sympleksu mniejszego do
większego, a parowaną od większego do mniejszego.
Wtedy $\mathcal V$ jest acykliczne wtedy i tylko wtedy,
gdy tak otrzymany graf skierowany nie ma cyklu skierowanego.
\end{stwierdzenie}

\begin{proof}
Zamknięta $\mathcal V$-ścieżka daje zamknięty spacer
skierowany: od $\beta_i$ przechodzimy strzałką parowaną
do $\alpha_i$, następnie strzałką nieparowaną do
$\beta_{i-1}$ (indeksy czytamy cyklicznie).
Każdy niepusty zamknięty spacer skierowany zawiera
cykl skierowany: wybieramy najkrótszy niepusty
zamknięty fragment spaceru; żaden jego wewnętrzny
wierzchołek nie może się powtarzać.
W drugą stronę niech $C$ będzie cyklem skierowanym
w grafie Hassego i wybierzmy na nim sympleks $\beta$
o maksymalnym wymiarze $p+1$. Strzałka wychodząca
z $\beta$ nie może prowadzić do nadsympleksu,
bo przeczyłoby to maksymalności wymiaru. Prowadzi
więc do pewnej ściany $\alpha$ wymiaru $p$ po
krawędzi parowanej. Ponieważ $\alpha$ jest już
sparowana z $\beta$, jej następna strzałka nie może
prowadzić w dół: oznaczałoby to drugą parę zawierającą
$\alpha$. Musi zatem prowadzić po krawędzi
nieparowanej do innego sympleksu wymiaru $p+1$.
Powtarzając rozumowanie, przechodzimy naprzemiennie
między wymiarami $p+1$ i $p$ (bez możliwości przejścia
na większy wymiar), aż cykl się zamknie. Otrzymaliśmy
niebanalną zamkniętą $\mathcal V$-ścieżkę,
czytając cykl w odwrotnym kierunku.
\end{proof}

\begin{uwaga}[Dyskretna funkcja Morse'a]
Równoważny opis można podać przy pomocy funkcji
$F$ na zbiorze sympleksów, dla której każdy sympleks
ma co najwyżej jedną bezpośrednią ścianę o wartości
$\ge F(\sigma)$ i co najwyżej jeden bezpośredni
nadsympleks o wartości $\le F(\sigma)$.
Po zadaniu acyklicznego parowania taki $F$ konstruujemy,
numerując wierzchołki skierowanego grafu Hassego
w porządku topologicznym (wzdłuż każdej strzałki
numer ściśle rośnie). Dla krawędzi nieparowanej
wartości rosną od ściany do nadsympleksu, a dla
parowanej maleją. Oba warunki definicji są więc
spełnione, ponieważ sympleks uczestniczy w co
najwyżej jednej parze.

Odwrotnie, parujemy $\alpha\prec\beta$, gdy
$F(\beta)\le F(\alpha)$. Warunki definicji wykluczają
dwie pary, w których ten sam sympleks jest dolnym
lub górnym elementem. Pozostaje wykluczyć sytuację
$\tau\prec\sigma\prec\beta$ i
$F(\beta)\le F(\sigma)\le F(\tau)$.
W kompleksie symplicjalnym istnieje wtedy druga
ściana $\sigma'\ne\sigma$ pomiędzy $\tau$ i $\beta$:
do $\tau$ dodajemy drugi z dwóch wierzchołków
$\beta\setminus\tau$. Warunek dla $\tau$ wymusza
$F(\sigma')>F(\tau)$, a warunek dla $\beta$ wymusza
$F(\sigma')<F(\beta)$, co jest sprzeczne.
Wzdłuż skierowanych krawędzi Hassego wartości $F$
nie maleją, a na krawędziach nieparowanych ściśle
rosną. Każdy cykl skierowany musiałby zawierać
krawędź nieparowaną (parowane krawędzie prowadzą
w dół wymiaru), więc cykl nie istnieje.
Otrzymujemy dokładnie acykliczne parowanie.
W dalszym wykładzie pracujemy bezpośrednio z nim.
\end{uwaga}

\begin{samepage}
\begin{center}
\begin{minipage}{.96\linewidth}\centering
\begin{tikzpicture}[scale=.91,>=Stealth,line cap=round]
 \node[tgTealVertex] (v0) at (0,0) {$a$};
 \node[tgBlueVertex] (v1) at (2,0) {$b$};
 \node[tgBlueVertex] (v2) at (4,0) {$c$};
 \draw[tgTree,->] (v1)--(v0);
 \draw[tgTree,->] (v2)--(v1);
 \node[text=tgNavy,font=\small] at (2,-.63)
   {pary: $(b,ab)$ i $(c,bc)$};
 \node[text=tgTeal,font=\small] at (0,.52) {krytyczny};
\end{tikzpicture}\par\smallskip

{\footnotesize Dla drzewa wybór rodzica każdego niekorzeniowego
wierzchołka daje acykliczne parowanie.}
\end{minipage}
\end{center}
\end{samepage}

\section{Dyskretny gradient na grafie i drzewa rozpinające}
\label{sec:tg-forman-on-graph}

\begin{twierdzenie}[Parowanie wyznaczone przez drzewo]
\label{thm:tg-forman-tree}
Niech $G$ będzie skończonym spójnym grafem z $n\ge1$
wierzchołkami i $m$ krawędziami, a $T$ drzewem
rozpinającym $G$. Dopuszczamy pętle i krawędzie równoległe, a krawędzie
traktujemy jako osobne obiekty. Parujemy wyłącznie
krawędzie drzewa, które mają dwa różne końce.
Dla grafu prostego jest to parowanie symplicjalne.
Dla ogólnego grafu rozumiemy ścieżki przez incydencję
wierzchołków i krawędzi; niesparowane pętle nigdy
nie są krokiem takiej ścieżki. Jeżeli potrzebujemy
regularnego kompleksu symplicjalnego, dzielimy
każdą krawędź na trzy krawędzie z dwoma nowymi,
osobnymi wierzchołkami. Otrzymany graf prosty ma
$n+2m$ wierzchołków i $3m$ krawędzi oraz tę samą
realizację topologiczną, więc liczba
$3m-(n+2m)+1=m-n+1$ nie zmienia się.
Wybierzmy korzeń $r\in V(T)$ i dla każdego $v\ne r$
sparujmy $v$ z krawędzią łączącą $v$ z jego rodzicem
w drzewie $T$. Otrzymane parowanie jest acykliczne.
Ma dokładnie jeden krytyczny wierzchołek oraz
$m-n+1$ krytycznych krawędzi.
\end{twierdzenie}

\begin{proof}
Każdy wierzchołek $v\ne r$ występuje dokładnie raz,
a każda krawędź drzewa również dokładnie raz jako
para z jej dalszym od korzenia końcem. Krawędzie
spoza $T$ nie są sparowane i są krytyczne; ich liczba
to $m-(n-1)=m-n+1$. Korzeń jest jedynym krytycznym
wierzchołkiem.

Każda $\mathcal V$-ścieżka wymiaru zero biegnie
od wierzchołka do jego rodzica: jeśli
$(v,e_v)$ jest parą, drugim końcem $e_v$ jest
właśnie rodzic $v$. Odległość w drzewie od korzenia
maleje więc o jeden przy każdym kroku takiej ścieżki.
Żadna niebanalna $\mathcal V$-ścieżka nie może być
zamknięta. W wymiarze jeden nie ma par krawędź--trójkąt,
więc innych zamkniętych ścieżek również nie ma.
\end{proof}

\begin{przyklad}[Trójkątny cykl]
Dla grafu $C_3$ o wierzchołkach $a,b,c$ wybieramy
drzewo $T=\{ab,bc\}$ i korzeń $a$.
Pary to $(b,ab)$ oraz $(c,bc)$. Krytyczny pozostaje
wierzchołek $a$ i krawędź $ac$. Otrzymujemy więc
jedną krytyczną komórkę wymiaru zero i jedną
wymiaru jeden, co odpowiada homotopijnemu typowi $S^1$.
Jeżeli teraz dołączymy wypełniający trójkąt $abc$,
możemy dodatkowo sparować krawędź $ac$ z trójkątem.
Ścieżki wymiaru zero pozostają bez zmian, a jedyna
para wymiarów $1,2$ prowadzi z $ac$ do $ab$ lub $bc$;
żadna z tych krawędzi nie jest dolnym elementem
kolejnej pary wymiarów $1,2$. Parowanie jest zatem
acykliczne. Pozostaje tylko krytyczny wierzchołek $a$,
zgodnie z kontraktowalnością dysku.
\end{przyklad}

\section{Elementarny kolaps a zachowanie typu homotopijnego}
\label{sec:tg-forman-collapse}

\begin{definicja}[Para wolna i elementarny kolaps]
Niech $K$ będzie skończonym kompleksem symplicjalnym.
Para $\alpha\prec\beta$ nazywa się \emph{wolna},
gdy $\beta$ jest jedynym maksymalnym sympleksem $K$
zawierającym $\alpha$. Usunięcie obu sympleksów
$\alpha$ i $\beta$ (bez usuwania ich pozostałych
ścian) nazywamy \emph{elementarnym kolapsem}.
Warunek $\alpha\prec\beta$ oznacza kowymiar jeden.
Każdy sympleks zawierający $\alpha$ jest zawarty
w jedynym maksymalnym sympleksie $\beta$, więc jest
równy $\alpha$ lub $\beta$. Żaden z pozostałych
sympleksów nie ma usuwanej ściany, zatem
$K'=K\setminus\{\alpha,\beta\}$ jest podkompleksem
(na zbiorze swoich wierzchołków).
\end{definicja}

\begin{stwierdzenie}[Kolaps nie zmienia typu homotopijnego]
Jeżeli $K'$ powstaje ze skończonego kompleksu
symplicjalnego $K$ przez elementarny kolaps wolnej
pary $\alpha\prec\beta$, to włączenie
$|K'|\hookrightarrow|K|$ jest równoważnością homotopijną,
a nawet istnieje silna retrakcja deformacyjna
$|K|$ na $|K'|$.
\end{stwierdzenie}

\begin{proof}
Niech $\beta=[v_0,\ldots,v_p,w]$, a jej wolną ścianą
będzie $\alpha=[v_0,\ldots,v_p]$. Każdy punkt
$|\beta|$ ma współrzędne barycentryczne
$(t_0,\ldots,t_p,s)$ o sumie $1$.
Oznaczmy $m=\min\{t_0,\ldots,t_p\}$.
Dla $u\in[0,1]$ definiujemy na $|\beta|$ homotopię
\[
H_u(t_0,\ldots,t_p,s)=
(t_0-um,\ldots,t_p-um,\ s+(p+1)um).
\]
Współrzędne pozostają nieujemne i nadal sumują się
do $1$. Jeżeli punkt należy do sumy pozostałych ścian
$\beta$ (czyli pewne $t_i=0$), to $m=0$ i $H_u$
pozostawia go nieruchomym. Dla $u=1$ pewna
współrzędna $t_i-m$ jest zerem, więc obraz leży
w tej sumie ścian. Homotopia jest ciągła,
$H_0=\mathrm{id}$, a $H_1$ jest retrakcją.

Z własności nośników współrzędnych barycentrycznych
wynika, że $|\beta|\cap|K'|$ jest dokładnie sumą
ścian $|\beta|$ innych niż usuwana ściana i sam
sympleks, czyli zbiorem punktów, dla których pewne
$t_i=0$. Oba zbiory $|\beta|$ i $|K'|$ są domknięte
i pokrywają $|K|$. Na ich przecięciu $H_u$ jest
identycznością dla wszystkich $u$. Lemat o sklejaniu
zastosowany do ich iloczynów z $[0,1]$ pozwala
przedłużyć homotopię identycznością na $|K'|$. Tak otrzymujemy silną retrakcję
deformacyjną $|K|$ na $|K'|$.
\end{proof}

\begin{przyklad}[Dysk z jednym krytycznym wierzchołkiem]
Niech $K$ będzie pełnym trójkątem $abc$. Najpierw
usuwamy wolną parę $(ac,abc)$, pozostawiając dwa
odcinki $ab$ i $bc$. Następnie usuwamy
$(c,bc)$, potem $(b,ab)$. Otrzymujemy jednoelementowy
kompleks $\{a\}$. Każdy krok jest elementarnym
kolapsem. W ten sposób nie tylko obliczamy homologie,
ale wykazujemy kontraktowalność geometryczną całej
realizacji $|K|$.
\end{przyklad}

\section{Komórki krytyczne, kompleks Morse'a i nierówności}
\label{sec:tg-forman-inequalities}

\begin{uwaga}[Acykliczne parowanie a elementarny kolaps]
Nie każda para z acyklicznego parowania jest wolną parą
w wyjściowym kompleksie. Nie wolno więc usuwać wszystkich
sparowanych komórek w dowolnej kolejności jako kolapsów
bez wcześniejszego sprawdzenia warunku wolnej ściany.
Twierdzenie Formana konstruuje odpowiedni model CW,
a jego dowód wymaga uporządkowania par zgodnie z dyskretnym
przepływem gradientowym. Jest to istotnie mocniejszy wynik
niż elementarne twierdzenie o pojedynczym kolapsie.
\end{uwaga}

\begin{twierdzenie}[Forman; podstawowe twierdzenie]
\label{thm:tg-forman-main}
Niech $K$ będzie skończonym regularnym kompleksem CW.
Parowanie komórek rozumiemy analogicznie jak dla kompleksów
symplicjalnych: parujemy komórkę $\alpha^p$ z komórką
$\beta^{p+1}$, dla której $\alpha$ jest ścianą
w brzegu $\overline\beta$, i każda komórka występuje
w co najwyżej jednej parze. Acykliczność oznacza brak
cykli skierowanych w grafie Hassego, z odwróconymi
krawędziami sparowanymi, zgodnie ze
Stwierdzeniem~\ref{thm:tg-morse-acyclic}.
Niech $\mathcal V$ będzie takim acyklicznym parowaniem.
Niech $c_p$ oznacza liczbę niesparowanych komórek
wymiaru $p$. Wtedy $|K|$ jest homotopijnie równoważny
kompleksowi CW mającemu dokładnie $c_p$ komórek
wymiaru $p$ dla każdego $p$.
W szczególności istnieje skończony kompleks łańcuchowy
Morse'a nad dowolnym ciałem $\mathbb F$, którego grupy łańcuchów
w stopniu $p$ mają wymiar $c_p$, a homologie są
izomorficzne z homologiami $|K|$.
\end{twierdzenie}

\begin{uwaga}[Status twierdzenia Formana]
Jest to głęboki rezultat; nie wynika on wyłącznie z
istnienia parowania. Dowód wykorzystuje porządkowanie
acyklicznych par i odpowiednie przekształcenia komórkowe
(algebraiczną wersję można skonstruować przez eliminację
par w kompleksie łańcuchowym). W szczególności pary
występujące w polu gradientowym \emph{nie muszą być
od razu wolnymi parami} w pierwotnym kompleksie.
Powyżej udowodniono bezpośrednio tylko zachowanie
homotopii dla rzeczywistych elementarnych kolapsów.
Pełne twierdzenie Formana przyjmujemy tutaj z literatury
(R. Forman, \emph{A User's Guide to Discrete Morse Theory},
Séminaire Lotharingien de Combinatoire 48, B48c, 2002).
Poniżej niezależnie udowodnimy jego część algebraiczną
dla kompleksów symplicjalnych, bez konstruowania modelu CW.
\end{uwaga}

\begin{stwierdzenie}[Algebraiczna eliminacja acyklicznego parowania]
\label{thm:tg-forman-algebraic-cancellation}
Niech $K$ będzie skończonym kompleksem symplicjalnym,
$\mathcal V$ acyklicznym parowaniem jego sympleksów,
a $\mathbb F$ dowolnym ciałem. Jeżeli $c_p$ oznacza
liczbę niesparowanych $p$-sympleksów, to symplicjalny
kompleks łańcuchowy $C_*(K;\mathbb F)$ jest
łańcuchowo homotopijnie równoważny skończonemu
kompleksowi $(M_*,d_*)$, dla którego
$\dim_{\mathbb F}M_p=c_p$ dla każdego $p$.
\end{stwierdzenie}

\begin{proof}
Najpierw opiszemy eliminację pojedynczej pary bazowej.
Niech $a$ będzie wektorem bazowym $C_p$,
$b$ wektorem bazowym $C_{p+1}$, a współczynnik
$a$ w $\partial b$ niech wynosi
$\lambda\in\mathbb F\setminus\{0\}$.
W bazie $C_p$ zastępujemy $a$ przez
$a'=\lambda^{-1}\partial b$; jest to zmiana bazy,
gdyż współczynnik $a$ w $a'$ wynosi $1$.
Mamy też $\partial a'=0$.
Dla każdego pozostałego wektora bazowego $x$
w $C_{p+1}$, oznaczając przez $\mu_x$
współczynnik $a$ w $\partial x$, wprowadzamy
$x'=x-\mu_x\lambda^{-1}b$.
Wówczas $\partial x'$ nie ma składowej przy $a'$,
a $\partial b=\lambda a'$. Istotnie, pisząc
$\partial b=\lambda a+\sum_{y\ne a}B_y y$ oraz
$\partial x=\mu_x a+\sum_{y\ne a}D_{yx}y$,
otrzymujemy
\[
 \partial x'=\sum_{y\ne a}
 (D_{yx}-B_y\lambda^{-1}\mu_x)y.
\]
Jest to dokładnie aktualizacja przez dopełnienie
Schura dla pozostałych wierszy i kolumn.
Dla dowolnego $z\in C_{p+2}$ współczynnik
$b$ w $\partial z$ zapisanym w nowej bazie również
jest zerowy: jeżeli ten współczynnik wynosi $\nu$, to
współczynnik $a'$ w $\partial^2 z$ jest równy
$\lambda\nu$, ponieważ pozostałe nowe generatory
stopnia $p+1$ nie mają składnika $a'$ w brzegu.
Z $\partial^2z=0$ i odwracalności $\lambda$ wynika
$\nu=0$.
Zatem kompleks rozkłada się na sumę prostą
dwuwyrazowego acyklicznego kompleksu
$\mathbb F b\xrightarrow{\lambda}\mathbb F a'$
i kompleksu na pozostałych wektorach bazowych.
Usunięcie pierwszego składnika jest łańcuchową
równoważnością homotopijną: na nim
$h(a')=\lambda^{-1}b$, $h(b)=0$ spełnia
$\partial h+h\partial=\mathrm{id}$ (sprawdzamy
na obu generatorach). Przedłużamy $h$ przez zero
na drugim składniku. Jeżeli $i$ i $\pi$ są jego
włączeniem i rzutem, to $\pi i=\mathrm{id}$ oraz
$\mathrm{id}-i\pi=\partial h+h\partial$.
Są to wymagane odwzorowania i homotopia łańcuchowa.

Pozostaje wykazać, że można wykonać eliminacje
dla wszystkich par $\mathcal V$.
Ustalmy wymiar $p$ i ponumerujmy pary
$(\alpha_i^p,\beta_i^{p+1})$.
W macierzy $\partial_{p+1}$, ograniczonej
do wierszy $\alpha_i$ i kolumn $\beta_j$,
elementy przekątnej są równe $\pm1$.
Niezerowy element poza przekątną w wierszu $i$
i kolumnie $j$ oznacza $\alpha_i\prec\beta_j$,
$i\ne j$, czyli przejście $\mathcal V$-ścieżki
od $\alpha_j$ do $\alpha_i$.
Acykliczność $\mathcal V$ sprawia, że graf
zależności par w ustalonym wymiarze jest
acykliczny. Numerujemy pary w porządku topologicznym
dla strzałek $j\to i$; niezerowy element poza
przekątną ma więc $j<i$. Macierz par jest dolnotrójkątna
z przekątną złożoną z $\pm1$.

Eliminujemy pary kolejno od pierwszej. W wierszu
pierwszej pary wszystkie współczynniki pozostałych
kolumn par są zerowe. W formule aktualizacji Schura
dla pozostałej macierzy par mamy zatem $\mu_x=0$,
więc jej wszystkie elementy pozostają niezmienione.
Indukcyjnie następna macierz par jest dolnotrójkątna
i każdy kolejny element przekątnej pozostaje jednostką.

Sprawdźmy jeszcze wpływ eliminacji w stopniach
$p,p+1$ na pozostałe różniczki. W stopniu $p$
pozostałe generatory $y\ne a$ nie zmieniają się,
więc macierz $\partial_p$ na nich jest po prostu
ograniczeniem do pozostałych kolumn. W stopniu $p+1$
generatory $x'$ różnią się od $x$ tylko wielokrotnością
usuwanego $b$. Dla $\partial_{p+2}z$ współczynniki
przy $x'$ są równe dawnym współczynnikom przy $x$:
podstawienie $x=x'+\mu_x\lambda^{-1}b$ zmienia
wyłącznie współczynnik przy $b$, który w nowej
bazie znika z warunku $\partial^2z=0$.
Macierz $\partial_{p+2}$ jest więc ograniczeniem
do pozostałych wierszy. Wszystkie inne różniczki
pozostają bez zmian. Ponieważ każda komórka należy
do co najwyżej jednej pary, usunięte kolumny lub
wiersze nie należą do macierzy par w sąsiednich
wymiarach. Ich współczynniki nie zmieniają się.
Możemy zatem eliminować kolejno całe grupy par
według wymiarów, stosując w każdej grupie ustalony
porządek topologiczny.

Po skończenie wielu krokach zostają dokładnie
niesparowane sympleksy każdego wymiaru.
Każdy krok zachowuje łańcuchowy typ homotopijny,
więc otrzymany kompleks $(M_*,d_*)$
ma żądane wymiary oraz izomorficzne homologie.
\end{proof}

\begin{wniosek}[Redukcja Morse'a nad liczbami całkowitymi]
Dla skończonego kompleksu symplicjalnego $K$
i acyklicznego parowania $\mathcal V$ istnieje kompleks
wolnych grup abelowych $M_*$, łańcuchowo homotopijnie
równoważny $C_*(K;\mathbb Z)$, z jednym generatorem
za każdy niesparowany sympleks odpowiedniego wymiaru.
\end{wniosek}

\begin{proof}
W dowodzie Stwierdzenia~\ref{thm:tg-forman-algebraic-cancellation}
każdy element główny początkowej macierzy par jest równy $1$ lub $-1$.
Są to elementy odwracalne także w $\mathbb Z$.
Trójkątność macierzy par zachowuje te wartości na
przekątnej podczas kolejnych eliminacji. Wszystkie
opisane zmiany baz, eliminacje i homotopie są więc
określone nad liczbami całkowitymi.
\end{proof}

\begin{uwaga}[Współczynniki i torsja]
Redukcja całkowitoliczbowa zachowuje również torsję
homologii. Nie można jej na ogół odczytać z samych
liczb Bettiego nad $\mathbb R$. Na przykład w modelu
CW płaszczyzny rzutowej $\mathbb{RP}^2$ jest po jednej
komórce w wymiarach $0,1,2$, a operatory brzegowe
to $\partial_1=0$ i $\partial_2=2$.
Stąd $H_1(\mathbb{RP}^2;\mathbb Z)\cong\mathbb Z/2\mathbb Z$,
ale $H_1(\mathbb{RP}^2;\mathbb R)=0$.
Nad $\mathbb F_2$ oba operatory znikają, więc
liczby Bettiego to $(1,1,1)$. Charakterystyka Eulera
jest równa $1$ przy każdym wyborze ciała.
\end{uwaga}

\begin{twierdzenie}[Słabe i silne nierówności Morse'a]
\label{thm:tg-forman-ineq}
Niech $K$ będzie skończonym kompleksem symplicjalnym,
$\mathcal V$ acyklicznym parowaniem, a $c_p$ liczbą
krytycznych $p$-sympleksów. Niech
$b_p=\dim_{\mathbb F}H_p(|K|;\mathbb F)$
dla ustalonego ciała $\mathbb F$. Wtedy dla każdego
$p\ge0$ zachodzi $c_p\ge b_p$, a ponadto
\[
 \sum_{j=0}^p(-1)^{p-j}c_j
 \ge \sum_{j=0}^p(-1)^{p-j}b_j,
 \qquad
 \sum_{j\ge0}(-1)^jc_j=\chi(|K|).
\]
\end{twierdzenie}

\begin{proof}
Na mocy Stwierdzenia~\ref{thm:tg-forman-algebraic-cancellation}
otrzymujemy kompleks łańcuchowy Morse'a $(M_*,d_*)$
z $\dim M_p=c_p$ i homologiami o wymiarach $b_p$.
Nie korzystamy w tym dowodzie z topologicznej
części twierdzenia Formana.
Niech $r_p=\operatorname{rank}d_p$; przyjmujemy
$r_0=0$ i $r_p=0$ ponad wymiarem kompleksu.
Z twierdzenia o rzędzie oraz definicji homologii mamy
$b_p=(c_p-r_p)-r_{p+1}$, zatem
\[
 c_p=b_p+r_p+r_{p+1}.
\]
Od razu dostajemy $c_p\ge b_p$.
Po zsumowaniu naprzemiennym od zera do $p$
wszystkie składniki $r_j$ w środku teleskopowo
się znoszą i pozostaje $r_{p+1}\ge0$. Stąd silne
nierówności. Sumując po wszystkich wymiarach,
ponieważ $r_{d+1}=0$ dla $d=\dim K$ (dla pustego
kompleksu wszystkie grupy są zerowe), uzyskujemy
$\sum(-1)^pc_p=\sum(-1)^pb_p=\chi(|K|)$,
gdzie ostatnia równość jest wzorem Eulera--Poincarégo
udowodnionym w rozdziale \ref{ch:tg-triangulacje}.
\end{proof}

\begin{przyklad}[Sfera z czterech trójkątów]
Rozważmy triangulację brzegu czworościanu o wierzchołkach
$0,1,2,3$. Parujemy wierzchołki $1,2,3$
odpowiednio z krawędziami $01,02,03$ oraz krawędzie
$12,13,23$ z trójkątami $012,013,023$.
Niesparowane są dokładnie wierzchołek $0$ i trójkąt
$123$. Parowanie jest acykliczne: w wymiarze zero
wszystkie ścieżki kierują się w stronę korzenia $0$,
a w wymiarze jeden po wejściu w sparowany trójkąt
przechodzimy tylko do krawędzi incydentnych z $0$,
które nie są sparowane z trójkątami, więc ścieżka
nie może się zamknąć. Otrzymujemy
$(c_0,c_1,c_2)=(1,0,1)$, zgodne z
$H_0(S^2)=H_2(S^2)\cong\mathbb Z$ i $H_1(S^2)=0$.
\end{przyklad}

\begin{samepage}
\begin{center}
\begin{minipage}{.96\linewidth}\centering
\begin{tikzpicture}[scale=1.0,line join=round,line cap=round]
 \coordinate (a) at (-1.15,-.7);\coordinate (b) at (1.2,-.7);
 \coordinate (c) at (.3,1.35);\coordinate (d) at (-.2,.18);
 \fill[tgGold!20] (b)--(c)--(d)--cycle;
 \draw[tgEdge] (a)--(b)--(c)--cycle (a)--(d)--(b) (c)--(d);
 \node[tgTealVertex] at (a) {$0$};
 \node[tgBlueVertex] at (b) {$1$};
 \node[tgBlueVertex] at (c) {$2$};
 \node[tgBlueVertex] at (d) {$3$};
 \node[font=\small,text=tgTeal] at (-1.55,-1.25) {krytyczne $0$};
 \node[font=\small,text=tgGold] at (1.6,.6) {krytyczne $123$};
\end{tikzpicture}\par\smallskip

{\footnotesize Na brzegu czworościanu pozostają dwie komórki
krytyczne: wierzchołek i ściana.}
\end{minipage}
\end{center}
\end{samepage}

\begin{twierdzenie}[Konstrukcja drzewo--drzewo dualne]
\label{thm:tg-forman-tree-cotree}
Niech $K$ będzie skończoną triangulacją spójnej,
zamkniętej powierzchni $M$ i niech $f_i$ oznacza
liczbę jej $i$-sympleksów. Istnieje acykliczne
parowanie sympleksów $K$, dla którego liczby
komórek krytycznych wynoszą
\[
(c_0,c_1,c_2)=(1,\ 2-\chi(M),\ 1).
\]
Jeżeli $M$ jest orientowalna rodzaju $g$, to
$(c_0,c_1,c_2)=(1,2g,1)$ i parowanie jest
doskonałe nad każdym ciałem.
\end{twierdzenie}

\begin{proof}
Szkielet jednowymiarowy $K^{(1)}$ jest spójny:
gdyby miał kilka składowych, każdy sympleks
miałby wszystkie wierzchołki w jednej z nich,
a realizacje odpowiadających podkompleksów
rozbiłyby $M$ na rozłączne niepuste zbiory
otwarte i domknięte. Wybierzmy jego drzewo
rozpinające $T$ i korzeń $r$.
Jak w Twierdzeniu~\ref{thm:tg-forman-tree}, dla
każdego $v\ne r$ parujemy wierzchołek $v$
z krawędzią drzewa łączącą go z rodzicem.
Otrzymujemy $f_0-1$ par wymiarów $0,1$.

Rozważmy graf dualny $G^*$ z
Sekcji~\ref{sec:tg-dual-triangulation}.
Pokażemy, że po usunięciu z niego krawędzi
odpowiadających krawędziom $T$ otrzymamy
graf spójny. Przypuśćmy przeciwnie, że
istnieje niepusty właściwy zbiór $A$
trójkątów, stanowiący sumę składowych
tak otrzymanego grafu. Nad $\mathbb F_2$
utwórzmy łańcuch $u=\sum_{\sigma\in A}\sigma$.
Ponieważ każda krawędź pierwotnej triangulacji
należy do dokładnie dwóch trójkątów,
łańcuch $z=\partial_2u$ składa się dokładnie
z krawędzi oddzielających trójkąty $A$
od pozostałych. Wszystkie należą do $T$,
gdyż innych krawędzi dualnych nie usuwaliśmy.
Ponadto $z\ne0$: pierwotny graf dualny
jest spójny, więc ma krawędź łączącą
$A$ z jego dopełnieniem.
Jednak $\partial_1z=\partial_1\partial_2u=0$.
Otrzymalibyśmy niezerowy cykl jednowymiarowy
nad $\mathbb F_2$ podparty na drzewie $T$,
co jest niemożliwe: podgraf złożony z krawędzi
o niezerowym współczynniku w $z$ i ich końców
jest lasem zawierającym co najmniej jedną krawędź.
Składowa z krawędzią ma liść, a współczynnik tego
liścia w $\partial_1z$ nad $\mathbb F_2$ wynosi $1$,
sprzecznie z $\partial_1z=0$.
Sprzeczność dowodzi spójności.

W pozostałym grafie dualnym wybierzmy
drzewo rozpinające $T^*$, ukorzenione
w wybranym trójkącie $\sigma_0$.
Dla każdego trójkąta $\sigma\ne\sigma_0$
niech $e_\sigma$ będzie krawędzią
pierwotnej triangulacji odpowiadającą
dualnej krawędzi łączącej $\sigma$
z jego rodzicem w $T^*$. Parujemy
$e_\sigma$ z $\sigma$. Jest to $f_2-1$
różnych par wymiarów $1,2$.
Żadna krawędź $e_\sigma$ nie należy do
$T$, zatem nie występuje w poprzednich
parach. Otrzymaliśmy poprawne parowanie.

Sprawdzimy acykliczność oddzielnie dla
$\mathcal V$-ścieżek każdego wymiaru.
W wymiarze zero kolejne wierzchołki
zbliżają się do korzenia $r$, więc
ich odległość w $T$ ściśle maleje.
W wymiarze jeden przejście z
$e_\sigma\prec\sigma$ do innej
sparowanej krawędzi $e_\tau\prec\sigma$
oznacza, że $\tau$ jest dzieckiem $\sigma$
w drzewie dualnym $T^*$: krawędź
$e_\sigma$ prowadząca do rodzica
została wykluczona w definicji
$\mathcal V$-ścieżki. Wobec tego
odległość trójkątów od korzenia
$\sigma_0$ ściśle rośnie i także
nie może powstać ścieżka zamknięta.
Innych rodzajów par nie ma.

Krytyczne pozostają: korzeń $r$,
trójkąt $\sigma_0$ oraz
$f_1-(f_0-1)-(f_2-1)$ krawędzi.
Ta ostatnia liczba wynosi
$2-(f_0-f_1+f_2)=2-\chi(M)$.
Dla orientowalnej powierzchni rodzaju
$g$ mamy $\chi(M)=2-2g$, a liczby Bettiego
nad dowolnym ciałem to $(1,2g,1)$.
Otrzymane parowanie jest więc doskonałe.
\end{proof}

\begin{przyklad}[Dyskretna funkcja Morse'a na torusie]
Dla standardowej siatki okresowej $3\times3$
na torusie, w której każdy kwadrat podzielono
przekątną na dwa trójkąty, mamy
$(f_0,f_1,f_2)=(9,27,18)$.
Konstrukcja z poprzedniego twierdzenia
paruje osiem wierzchołków z krawędziami
drzewa rozpinającego i siedemnaście
dalszych krawędzi z trójkątami drzewa
dualnego. Pozostają dokładnie
$(c_0,c_1,c_2)=(1,2,1)$, zgodnie
z homologiami torusa.
\end{przyklad}

\section{Zestawienie teorii ciągłej i dyskretnej}
\label{sec:tg-morse-bridge}

Oba rodzaje teorii Morse'a prowadzą do opisu topologii
za pomocą niewielkiej liczby danych krytycznych,
ale nie są definicyjnie tożsame. Klasyczna funkcja
Morse'a na rozmaitości jest gładka, a jej
punkty krytyczne są badane przy użyciu hesjanu.
Dyskretne pole gradientowe działa na zbiorze komórek,
a warunkiem podstawowym jest brak zamkniętych
$\mathcal V$-ścieżek. W obu przypadkach otrzymujemy
nierówności wiążące liczby krytycznych obiektów
z liczbami Bettiego i charakterystyką Eulera.

\begin{uwaga}[Minimalność a doskonałość]
Parowanie nazywamy \emph{doskonałym nad ciałem}
$\mathbb F$, gdy $c_p=b_p$ dla wszystkich $p$.
Jest to warunek silniejszy niż samo spełnianie równości
$\sum(-1)^pc_p=\chi$. Dla grafu spójnego konstrukcja
z drzewa rozpinającego daje doskonałe parowanie
($c_0=1$, $c_1=m-n+1$). Dla brzegu czworościanu
otrzymujemy również doskonałe parowanie z dwiema
komórkami krytycznymi. W bardziej złożonych
kompleksach istnienie doskonałego parowania nie jest
zagwarantowane.
\end{uwaga}

